\documentclass[11pt]{article}

    \usepackage[breakable]{tcolorbox}
    \usepackage{parskip} % Stop auto-indenting (to mimic markdown behaviour)
    

    % Basic figure setup, for now with no caption control since it's done
    % automatically by Pandoc (which extracts ![](path) syntax from Markdown).
    \usepackage{graphicx}
    % Keep aspect ratio if custom image width or height is specified
    \setkeys{Gin}{keepaspectratio}
    % Maintain compatibility with old templates. Remove in nbconvert 6.0
    \let\Oldincludegraphics\includegraphics
    % Ensure that by default, figures have no caption (until we provide a
    % proper Figure object with a Caption API and a way to capture that
    % in the conversion process - todo).
    \usepackage{caption}
    \DeclareCaptionFormat{nocaption}{}
    \captionsetup{format=nocaption,aboveskip=0pt,belowskip=0pt}

    \usepackage{float}
    \floatplacement{figure}{H} % forces figures to be placed at the correct location
    \usepackage{xcolor} % Allow colors to be defined
    \usepackage{enumerate} % Needed for markdown enumerations to work
    \usepackage{geometry} % Used to adjust the document margins
    \usepackage{amsmath} % Equations
    \usepackage{amssymb} % Equations
    \usepackage{textcomp} % defines textquotesingle
    % Hack from http://tex.stackexchange.com/a/47451/13684:
    \AtBeginDocument{%
        \def\PYZsq{\textquotesingle}% Upright quotes in Pygmentized code
    }
    \usepackage{upquote} % Upright quotes for verbatim code
    \usepackage{eurosym} % defines \euro

    \usepackage{iftex}
    \ifPDFTeX
        \usepackage[T1]{fontenc}
        \IfFileExists{alphabeta.sty}{
              \usepackage{alphabeta}
          }{
              \usepackage[mathletters]{ucs}
              \usepackage[utf8x]{inputenc}
          }
    \else
        \usepackage{fontspec}
        \usepackage{unicode-math}
    \fi

    \usepackage{fancyvrb} % verbatim replacement that allows latex
    \usepackage{grffile} % extends the file name processing of package graphics
                         % to support a larger range
    \makeatletter % fix for old versions of grffile with XeLaTeX
    \@ifpackagelater{grffile}{2019/11/01}
    {
      % Do nothing on new versions
    }
    {
      \def\Gread@@xetex#1{%
        \IfFileExists{"\Gin@base".bb}%
        {\Gread@eps{\Gin@base.bb}}%
        {\Gread@@xetex@aux#1}%
      }
    }
    \makeatother
    \usepackage[Export]{adjustbox} % Used to constrain images to a maximum size
    \adjustboxset{max size={0.9\linewidth}{0.9\paperheight}}

    % The hyperref package gives us a pdf with properly built
    % internal navigation ('pdf bookmarks' for the table of contents,
    % internal cross-reference links, web links for URLs, etc.)
    \usepackage{hyperref}
    % The default LaTeX title has an obnoxious amount of whitespace. By default,
    % titling removes some of it. It also provides customization options.
    \usepackage{titling}
    \usepackage{longtable} % longtable support required by pandoc >1.10
    \usepackage{booktabs}  % table support for pandoc > 1.12.2
    \usepackage{array}     % table support for pandoc >= 2.11.3
    \usepackage{calc}      % table minipage width calculation for pandoc >= 2.11.1
    \usepackage[inline]{enumitem} % IRkernel/repr support (it uses the enumerate* environment)
    \usepackage[normalem]{ulem} % ulem is needed to support strikethroughs (\sout)
                                % normalem makes italics be italics, not underlines
    \usepackage{soul}      % strikethrough (\st) support for pandoc >= 3.0.0
    \usepackage{mathrsfs}
    

    
    % Colors for the hyperref package
    \definecolor{urlcolor}{rgb}{0,.145,.698}
    \definecolor{linkcolor}{rgb}{.71,0.21,0.01}
    \definecolor{citecolor}{rgb}{.12,.54,.11}

    % ANSI colors
    \definecolor{ansi-black}{HTML}{3E424D}
    \definecolor{ansi-black-intense}{HTML}{282C36}
    \definecolor{ansi-red}{HTML}{E75C58}
    \definecolor{ansi-red-intense}{HTML}{B22B31}
    \definecolor{ansi-green}{HTML}{00A250}
    \definecolor{ansi-green-intense}{HTML}{007427}
    \definecolor{ansi-yellow}{HTML}{DDB62B}
    \definecolor{ansi-yellow-intense}{HTML}{B27D12}
    \definecolor{ansi-blue}{HTML}{208FFB}
    \definecolor{ansi-blue-intense}{HTML}{0065CA}
    \definecolor{ansi-magenta}{HTML}{D160C4}
    \definecolor{ansi-magenta-intense}{HTML}{A03196}
    \definecolor{ansi-cyan}{HTML}{60C6C8}
    \definecolor{ansi-cyan-intense}{HTML}{258F8F}
    \definecolor{ansi-white}{HTML}{C5C1B4}
    \definecolor{ansi-white-intense}{HTML}{A1A6B2}
    \definecolor{ansi-default-inverse-fg}{HTML}{FFFFFF}
    \definecolor{ansi-default-inverse-bg}{HTML}{000000}

    % common color for the border for error outputs.
    \definecolor{outerrorbackground}{HTML}{FFDFDF}

    % commands and environments needed by pandoc snippets
    % extracted from the output of `pandoc -s`
    \providecommand{\tightlist}{%
      \setlength{\itemsep}{0pt}\setlength{\parskip}{0pt}}
    \DefineVerbatimEnvironment{Highlighting}{Verbatim}{commandchars=\\\{\}}
    % Add ',fontsize=\small' for more characters per line
    \newenvironment{Shaded}{}{}
    \newcommand{\KeywordTok}[1]{\textcolor[rgb]{0.00,0.44,0.13}{\textbf{{#1}}}}
    \newcommand{\DataTypeTok}[1]{\textcolor[rgb]{0.56,0.13,0.00}{{#1}}}
    \newcommand{\DecValTok}[1]{\textcolor[rgb]{0.25,0.63,0.44}{{#1}}}
    \newcommand{\BaseNTok}[1]{\textcolor[rgb]{0.25,0.63,0.44}{{#1}}}
    \newcommand{\FloatTok}[1]{\textcolor[rgb]{0.25,0.63,0.44}{{#1}}}
    \newcommand{\CharTok}[1]{\textcolor[rgb]{0.25,0.44,0.63}{{#1}}}
    \newcommand{\StringTok}[1]{\textcolor[rgb]{0.25,0.44,0.63}{{#1}}}
    \newcommand{\CommentTok}[1]{\textcolor[rgb]{0.38,0.63,0.69}{\textit{{#1}}}}
    \newcommand{\OtherTok}[1]{\textcolor[rgb]{0.00,0.44,0.13}{{#1}}}
    \newcommand{\AlertTok}[1]{\textcolor[rgb]{1.00,0.00,0.00}{\textbf{{#1}}}}
    \newcommand{\FunctionTok}[1]{\textcolor[rgb]{0.02,0.16,0.49}{{#1}}}
    \newcommand{\RegionMarkerTok}[1]{{#1}}
    \newcommand{\ErrorTok}[1]{\textcolor[rgb]{1.00,0.00,0.00}{\textbf{{#1}}}}
    \newcommand{\NormalTok}[1]{{#1}}

    % Additional commands for more recent versions of Pandoc
    \newcommand{\ConstantTok}[1]{\textcolor[rgb]{0.53,0.00,0.00}{{#1}}}
    \newcommand{\SpecialCharTok}[1]{\textcolor[rgb]{0.25,0.44,0.63}{{#1}}}
    \newcommand{\VerbatimStringTok}[1]{\textcolor[rgb]{0.25,0.44,0.63}{{#1}}}
    \newcommand{\SpecialStringTok}[1]{\textcolor[rgb]{0.73,0.40,0.53}{{#1}}}
    \newcommand{\ImportTok}[1]{{#1}}
    \newcommand{\DocumentationTok}[1]{\textcolor[rgb]{0.73,0.13,0.13}{\textit{{#1}}}}
    \newcommand{\AnnotationTok}[1]{\textcolor[rgb]{0.38,0.63,0.69}{\textbf{\textit{{#1}}}}}
    \newcommand{\CommentVarTok}[1]{\textcolor[rgb]{0.38,0.63,0.69}{\textbf{\textit{{#1}}}}}
    \newcommand{\VariableTok}[1]{\textcolor[rgb]{0.10,0.09,0.49}{{#1}}}
    \newcommand{\ControlFlowTok}[1]{\textcolor[rgb]{0.00,0.44,0.13}{\textbf{{#1}}}}
    \newcommand{\OperatorTok}[1]{\textcolor[rgb]{0.40,0.40,0.40}{{#1}}}
    \newcommand{\BuiltInTok}[1]{{#1}}
    \newcommand{\ExtensionTok}[1]{{#1}}
    \newcommand{\PreprocessorTok}[1]{\textcolor[rgb]{0.74,0.48,0.00}{{#1}}}
    \newcommand{\AttributeTok}[1]{\textcolor[rgb]{0.49,0.56,0.16}{{#1}}}
    \newcommand{\InformationTok}[1]{\textcolor[rgb]{0.38,0.63,0.69}{\textbf{\textit{{#1}}}}}
    \newcommand{\WarningTok}[1]{\textcolor[rgb]{0.38,0.63,0.69}{\textbf{\textit{{#1}}}}}
    \makeatletter
    \newsavebox\pandoc@box
    \newcommand*\pandocbounded[1]{%
      \sbox\pandoc@box{#1}%
      % scaling factors for width and height
      \Gscale@div\@tempa\textheight{\dimexpr\ht\pandoc@box+\dp\pandoc@box\relax}%
      \Gscale@div\@tempb\linewidth{\wd\pandoc@box}%
      % select the smaller of both
      \ifdim\@tempb\p@<\@tempa\p@
        \let\@tempa\@tempb
      \fi
      % scaling accordingly (\@tempa < 1)
      \ifdim\@tempa\p@<\p@
        \scalebox{\@tempa}{\usebox\pandoc@box}%
      % scaling not needed, use as it is
      \else
        \usebox{\pandoc@box}%
      \fi
    }
    \makeatother

    % Define a nice break command that doesn't care if a line doesn't already
    % exist.
    \def\br{\hspace*{\fill} \\* }
    % Math Jax compatibility definitions
    \def\gt{>}
    \def\lt{<}
    \let\Oldtex\TeX
    \let\Oldlatex\LaTeX
    \renewcommand{\TeX}{\textrm{\Oldtex}}
    \renewcommand{\LaTeX}{\textrm{\Oldlatex}}
    % Document parameters
    % Document title
    \title{InterfazPaqNOSup}
    
    
    
    
    
    
    
% Pygments definitions
\makeatletter
\def\PY@reset{\let\PY@it=\relax \let\PY@bf=\relax%
    \let\PY@ul=\relax \let\PY@tc=\relax%
    \let\PY@bc=\relax \let\PY@ff=\relax}
\def\PY@tok#1{\csname PY@tok@#1\endcsname}
\def\PY@toks#1+{\ifx\relax#1\empty\else%
    \PY@tok{#1}\expandafter\PY@toks\fi}
\def\PY@do#1{\PY@bc{\PY@tc{\PY@ul{%
    \PY@it{\PY@bf{\PY@ff{#1}}}}}}}
\def\PY#1#2{\PY@reset\PY@toks#1+\relax+\PY@do{#2}}

\@namedef{PY@tok@w}{\def\PY@tc##1{\textcolor[rgb]{0.73,0.73,0.73}{##1}}}
\@namedef{PY@tok@c}{\let\PY@it=\textit\def\PY@tc##1{\textcolor[rgb]{0.24,0.48,0.48}{##1}}}
\@namedef{PY@tok@cp}{\def\PY@tc##1{\textcolor[rgb]{0.61,0.40,0.00}{##1}}}
\@namedef{PY@tok@k}{\let\PY@bf=\textbf\def\PY@tc##1{\textcolor[rgb]{0.00,0.50,0.00}{##1}}}
\@namedef{PY@tok@kp}{\def\PY@tc##1{\textcolor[rgb]{0.00,0.50,0.00}{##1}}}
\@namedef{PY@tok@kt}{\def\PY@tc##1{\textcolor[rgb]{0.69,0.00,0.25}{##1}}}
\@namedef{PY@tok@o}{\def\PY@tc##1{\textcolor[rgb]{0.40,0.40,0.40}{##1}}}
\@namedef{PY@tok@ow}{\let\PY@bf=\textbf\def\PY@tc##1{\textcolor[rgb]{0.67,0.13,1.00}{##1}}}
\@namedef{PY@tok@nb}{\def\PY@tc##1{\textcolor[rgb]{0.00,0.50,0.00}{##1}}}
\@namedef{PY@tok@nf}{\def\PY@tc##1{\textcolor[rgb]{0.00,0.00,1.00}{##1}}}
\@namedef{PY@tok@nc}{\let\PY@bf=\textbf\def\PY@tc##1{\textcolor[rgb]{0.00,0.00,1.00}{##1}}}
\@namedef{PY@tok@nn}{\let\PY@bf=\textbf\def\PY@tc##1{\textcolor[rgb]{0.00,0.00,1.00}{##1}}}
\@namedef{PY@tok@ne}{\let\PY@bf=\textbf\def\PY@tc##1{\textcolor[rgb]{0.80,0.25,0.22}{##1}}}
\@namedef{PY@tok@nv}{\def\PY@tc##1{\textcolor[rgb]{0.10,0.09,0.49}{##1}}}
\@namedef{PY@tok@no}{\def\PY@tc##1{\textcolor[rgb]{0.53,0.00,0.00}{##1}}}
\@namedef{PY@tok@nl}{\def\PY@tc##1{\textcolor[rgb]{0.46,0.46,0.00}{##1}}}
\@namedef{PY@tok@ni}{\let\PY@bf=\textbf\def\PY@tc##1{\textcolor[rgb]{0.44,0.44,0.44}{##1}}}
\@namedef{PY@tok@na}{\def\PY@tc##1{\textcolor[rgb]{0.41,0.47,0.13}{##1}}}
\@namedef{PY@tok@nt}{\let\PY@bf=\textbf\def\PY@tc##1{\textcolor[rgb]{0.00,0.50,0.00}{##1}}}
\@namedef{PY@tok@nd}{\def\PY@tc##1{\textcolor[rgb]{0.67,0.13,1.00}{##1}}}
\@namedef{PY@tok@s}{\def\PY@tc##1{\textcolor[rgb]{0.73,0.13,0.13}{##1}}}
\@namedef{PY@tok@sd}{\let\PY@it=\textit\def\PY@tc##1{\textcolor[rgb]{0.73,0.13,0.13}{##1}}}
\@namedef{PY@tok@si}{\let\PY@bf=\textbf\def\PY@tc##1{\textcolor[rgb]{0.64,0.35,0.47}{##1}}}
\@namedef{PY@tok@se}{\let\PY@bf=\textbf\def\PY@tc##1{\textcolor[rgb]{0.67,0.36,0.12}{##1}}}
\@namedef{PY@tok@sr}{\def\PY@tc##1{\textcolor[rgb]{0.64,0.35,0.47}{##1}}}
\@namedef{PY@tok@ss}{\def\PY@tc##1{\textcolor[rgb]{0.10,0.09,0.49}{##1}}}
\@namedef{PY@tok@sx}{\def\PY@tc##1{\textcolor[rgb]{0.00,0.50,0.00}{##1}}}
\@namedef{PY@tok@m}{\def\PY@tc##1{\textcolor[rgb]{0.40,0.40,0.40}{##1}}}
\@namedef{PY@tok@gh}{\let\PY@bf=\textbf\def\PY@tc##1{\textcolor[rgb]{0.00,0.00,0.50}{##1}}}
\@namedef{PY@tok@gu}{\let\PY@bf=\textbf\def\PY@tc##1{\textcolor[rgb]{0.50,0.00,0.50}{##1}}}
\@namedef{PY@tok@gd}{\def\PY@tc##1{\textcolor[rgb]{0.63,0.00,0.00}{##1}}}
\@namedef{PY@tok@gi}{\def\PY@tc##1{\textcolor[rgb]{0.00,0.52,0.00}{##1}}}
\@namedef{PY@tok@gr}{\def\PY@tc##1{\textcolor[rgb]{0.89,0.00,0.00}{##1}}}
\@namedef{PY@tok@ge}{\let\PY@it=\textit}
\@namedef{PY@tok@gs}{\let\PY@bf=\textbf}
\@namedef{PY@tok@ges}{\let\PY@bf=\textbf\let\PY@it=\textit}
\@namedef{PY@tok@gp}{\let\PY@bf=\textbf\def\PY@tc##1{\textcolor[rgb]{0.00,0.00,0.50}{##1}}}
\@namedef{PY@tok@go}{\def\PY@tc##1{\textcolor[rgb]{0.44,0.44,0.44}{##1}}}
\@namedef{PY@tok@gt}{\def\PY@tc##1{\textcolor[rgb]{0.00,0.27,0.87}{##1}}}
\@namedef{PY@tok@err}{\def\PY@bc##1{{\setlength{\fboxsep}{\string -\fboxrule}\fcolorbox[rgb]{1.00,0.00,0.00}{1,1,1}{\strut ##1}}}}
\@namedef{PY@tok@kc}{\let\PY@bf=\textbf\def\PY@tc##1{\textcolor[rgb]{0.00,0.50,0.00}{##1}}}
\@namedef{PY@tok@kd}{\let\PY@bf=\textbf\def\PY@tc##1{\textcolor[rgb]{0.00,0.50,0.00}{##1}}}
\@namedef{PY@tok@kn}{\let\PY@bf=\textbf\def\PY@tc##1{\textcolor[rgb]{0.00,0.50,0.00}{##1}}}
\@namedef{PY@tok@kr}{\let\PY@bf=\textbf\def\PY@tc##1{\textcolor[rgb]{0.00,0.50,0.00}{##1}}}
\@namedef{PY@tok@bp}{\def\PY@tc##1{\textcolor[rgb]{0.00,0.50,0.00}{##1}}}
\@namedef{PY@tok@fm}{\def\PY@tc##1{\textcolor[rgb]{0.00,0.00,1.00}{##1}}}
\@namedef{PY@tok@vc}{\def\PY@tc##1{\textcolor[rgb]{0.10,0.09,0.49}{##1}}}
\@namedef{PY@tok@vg}{\def\PY@tc##1{\textcolor[rgb]{0.10,0.09,0.49}{##1}}}
\@namedef{PY@tok@vi}{\def\PY@tc##1{\textcolor[rgb]{0.10,0.09,0.49}{##1}}}
\@namedef{PY@tok@vm}{\def\PY@tc##1{\textcolor[rgb]{0.10,0.09,0.49}{##1}}}
\@namedef{PY@tok@sa}{\def\PY@tc##1{\textcolor[rgb]{0.73,0.13,0.13}{##1}}}
\@namedef{PY@tok@sb}{\def\PY@tc##1{\textcolor[rgb]{0.73,0.13,0.13}{##1}}}
\@namedef{PY@tok@sc}{\def\PY@tc##1{\textcolor[rgb]{0.73,0.13,0.13}{##1}}}
\@namedef{PY@tok@dl}{\def\PY@tc##1{\textcolor[rgb]{0.73,0.13,0.13}{##1}}}
\@namedef{PY@tok@s2}{\def\PY@tc##1{\textcolor[rgb]{0.73,0.13,0.13}{##1}}}
\@namedef{PY@tok@sh}{\def\PY@tc##1{\textcolor[rgb]{0.73,0.13,0.13}{##1}}}
\@namedef{PY@tok@s1}{\def\PY@tc##1{\textcolor[rgb]{0.73,0.13,0.13}{##1}}}
\@namedef{PY@tok@mb}{\def\PY@tc##1{\textcolor[rgb]{0.40,0.40,0.40}{##1}}}
\@namedef{PY@tok@mf}{\def\PY@tc##1{\textcolor[rgb]{0.40,0.40,0.40}{##1}}}
\@namedef{PY@tok@mh}{\def\PY@tc##1{\textcolor[rgb]{0.40,0.40,0.40}{##1}}}
\@namedef{PY@tok@mi}{\def\PY@tc##1{\textcolor[rgb]{0.40,0.40,0.40}{##1}}}
\@namedef{PY@tok@il}{\def\PY@tc##1{\textcolor[rgb]{0.40,0.40,0.40}{##1}}}
\@namedef{PY@tok@mo}{\def\PY@tc##1{\textcolor[rgb]{0.40,0.40,0.40}{##1}}}
\@namedef{PY@tok@ch}{\let\PY@it=\textit\def\PY@tc##1{\textcolor[rgb]{0.24,0.48,0.48}{##1}}}
\@namedef{PY@tok@cm}{\let\PY@it=\textit\def\PY@tc##1{\textcolor[rgb]{0.24,0.48,0.48}{##1}}}
\@namedef{PY@tok@cpf}{\let\PY@it=\textit\def\PY@tc##1{\textcolor[rgb]{0.24,0.48,0.48}{##1}}}
\@namedef{PY@tok@c1}{\let\PY@it=\textit\def\PY@tc##1{\textcolor[rgb]{0.24,0.48,0.48}{##1}}}
\@namedef{PY@tok@cs}{\let\PY@it=\textit\def\PY@tc##1{\textcolor[rgb]{0.24,0.48,0.48}{##1}}}

\def\PYZbs{\char`\\}
\def\PYZus{\char`\_}
\def\PYZob{\char`\{}
\def\PYZcb{\char`\}}
\def\PYZca{\char`\^}
\def\PYZam{\char`\&}
\def\PYZlt{\char`\<}
\def\PYZgt{\char`\>}
\def\PYZsh{\char`\#}
\def\PYZpc{\char`\%}
\def\PYZdl{\char`\$}
\def\PYZhy{\char`\-}
\def\PYZsq{\char`\'}
\def\PYZdq{\char`\"}
\def\PYZti{\char`\~}
% for compatibility with earlier versions
\def\PYZat{@}
\def\PYZlb{[}
\def\PYZrb{]}
\makeatother


    % For linebreaks inside Verbatim environment from package fancyvrb.
    \makeatletter
        \newbox\Wrappedcontinuationbox
        \newbox\Wrappedvisiblespacebox
        \newcommand*\Wrappedvisiblespace {\textcolor{red}{\textvisiblespace}}
        \newcommand*\Wrappedcontinuationsymbol {\textcolor{red}{\llap{\tiny$\m@th\hookrightarrow$}}}
        \newcommand*\Wrappedcontinuationindent {3ex }
        \newcommand*\Wrappedafterbreak {\kern\Wrappedcontinuationindent\copy\Wrappedcontinuationbox}
        % Take advantage of the already applied Pygments mark-up to insert
        % potential linebreaks for TeX processing.
        %        {, <, #, %, $, ' and ": go to next line.
        %        _, }, ^, &, >, - and ~: stay at end of broken line.
        % Use of \textquotesingle for straight quote.
        \newcommand*\Wrappedbreaksatspecials {%
            \def\PYGZus{\discretionary{\char`\_}{\Wrappedafterbreak}{\char`\_}}%
            \def\PYGZob{\discretionary{}{\Wrappedafterbreak\char`\{}{\char`\{}}%
            \def\PYGZcb{\discretionary{\char`\}}{\Wrappedafterbreak}{\char`\}}}%
            \def\PYGZca{\discretionary{\char`\^}{\Wrappedafterbreak}{\char`\^}}%
            \def\PYGZam{\discretionary{\char`\&}{\Wrappedafterbreak}{\char`\&}}%
            \def\PYGZlt{\discretionary{}{\Wrappedafterbreak\char`\<}{\char`\<}}%
            \def\PYGZgt{\discretionary{\char`\>}{\Wrappedafterbreak}{\char`\>}}%
            \def\PYGZsh{\discretionary{}{\Wrappedafterbreak\char`\#}{\char`\#}}%
            \def\PYGZpc{\discretionary{}{\Wrappedafterbreak\char`\%}{\char`\%}}%
            \def\PYGZdl{\discretionary{}{\Wrappedafterbreak\char`\$}{\char`\$}}%
            \def\PYGZhy{\discretionary{\char`\-}{\Wrappedafterbreak}{\char`\-}}%
            \def\PYGZsq{\discretionary{}{\Wrappedafterbreak\textquotesingle}{\textquotesingle}}%
            \def\PYGZdq{\discretionary{}{\Wrappedafterbreak\char`\"}{\char`\"}}%
            \def\PYGZti{\discretionary{\char`\~}{\Wrappedafterbreak}{\char`\~}}%
        }
        % Some characters . , ; ? ! / are not pygmentized.
        % This macro makes them "active" and they will insert potential linebreaks
        \newcommand*\Wrappedbreaksatpunct {%
            \lccode`\~`\.\lowercase{\def~}{\discretionary{\hbox{\char`\.}}{\Wrappedafterbreak}{\hbox{\char`\.}}}%
            \lccode`\~`\,\lowercase{\def~}{\discretionary{\hbox{\char`\,}}{\Wrappedafterbreak}{\hbox{\char`\,}}}%
            \lccode`\~`\;\lowercase{\def~}{\discretionary{\hbox{\char`\;}}{\Wrappedafterbreak}{\hbox{\char`\;}}}%
            \lccode`\~`\:\lowercase{\def~}{\discretionary{\hbox{\char`\:}}{\Wrappedafterbreak}{\hbox{\char`\:}}}%
            \lccode`\~`\?\lowercase{\def~}{\discretionary{\hbox{\char`\?}}{\Wrappedafterbreak}{\hbox{\char`\?}}}%
            \lccode`\~`\!\lowercase{\def~}{\discretionary{\hbox{\char`\!}}{\Wrappedafterbreak}{\hbox{\char`\!}}}%
            \lccode`\~`\/\lowercase{\def~}{\discretionary{\hbox{\char`\/}}{\Wrappedafterbreak}{\hbox{\char`\/}}}%
            \catcode`\.\active
            \catcode`\,\active
            \catcode`\;\active
            \catcode`\:\active
            \catcode`\?\active
            \catcode`\!\active
            \catcode`\/\active
            \lccode`\~`\~
        }
    \makeatother

    \let\OriginalVerbatim=\Verbatim
    \makeatletter
    \renewcommand{\Verbatim}[1][1]{%
        %\parskip\z@skip
        \sbox\Wrappedcontinuationbox {\Wrappedcontinuationsymbol}%
        \sbox\Wrappedvisiblespacebox {\FV@SetupFont\Wrappedvisiblespace}%
        \def\FancyVerbFormatLine ##1{\hsize\linewidth
            \vtop{\raggedright\hyphenpenalty\z@\exhyphenpenalty\z@
                \doublehyphendemerits\z@\finalhyphendemerits\z@
                \strut ##1\strut}%
        }%
        % If the linebreak is at a space, the latter will be displayed as visible
        % space at end of first line, and a continuation symbol starts next line.
        % Stretch/shrink are however usually zero for typewriter font.
        \def\FV@Space {%
            \nobreak\hskip\z@ plus\fontdimen3\font minus\fontdimen4\font
            \discretionary{\copy\Wrappedvisiblespacebox}{\Wrappedafterbreak}
            {\kern\fontdimen2\font}%
        }%

        % Allow breaks at special characters using \PYG... macros.
        \Wrappedbreaksatspecials
        % Breaks at punctuation characters . , ; ? ! and / need catcode=\active
        \OriginalVerbatim[#1,codes*=\Wrappedbreaksatpunct]%
    }
    \makeatother

    % Exact colors from NB
    \definecolor{incolor}{HTML}{303F9F}
    \definecolor{outcolor}{HTML}{D84315}
    \definecolor{cellborder}{HTML}{CFCFCF}
    \definecolor{cellbackground}{HTML}{F7F7F7}

    % prompt
    \makeatletter
    \newcommand{\boxspacing}{\kern\kvtcb@left@rule\kern\kvtcb@boxsep}
    \makeatother
    \newcommand{\prompt}[4]{
        {\ttfamily\llap{{\color{#2}[#3]:\hspace{3pt}#4}}\vspace{-\baselineskip}}
    }
    

    
    % Prevent overflowing lines due to hard-to-break entities
    \sloppy
    % Setup hyperref package
    \hypersetup{
      breaklinks=true,  % so long urls are correctly broken across lines
      colorlinks=true,
      urlcolor=urlcolor,
      linkcolor=linkcolor,
      citecolor=citecolor,
      }
    % Slightly bigger margins than the latex defaults
    
    \geometry{verbose,tmargin=1in,bmargin=1in,lmargin=1in,rmargin=1in}
    
    

\begin{document}
    
    \maketitle
    
    

    
    \section{Caso \#1}\label{caso-1}

\subsubsection{Desarrollo de Clusters}\label{desarrollo-de-clusters}

Dr.~Juan de Dios Murillo-Morera

    

    

    \subsection{Ejecucion de los clusters usando el paquete creado
M\_Caso1y2.py}\label{ejecucion-de-los-clusters-usando-el-paquete-creado-m_caso1y2.py}

    \begin{tcolorbox}[breakable, size=fbox, boxrule=1pt, pad at break*=1mm,colback=cellbackground, colframe=cellborder]
\prompt{In}{incolor}{1}{\boxspacing}
\begin{Verbatim}[commandchars=\\\{\}]
\PY{k+kn}{import}\PY{+w}{ }\PY{n+nn}{PaqNOSup}\PY{+w}{ }\PY{k}{as}\PY{+w}{ }\PY{n+nn}{noSup}
\end{Verbatim}
\end{tcolorbox}

    \subsubsection{Primer Dataset: Hotel}\label{primer-dataset-hotel}

    \begin{tcolorbox}[breakable, size=fbox, boxrule=1pt, pad at break*=1mm,colback=cellbackground, colframe=cellborder]
\prompt{In}{incolor}{2}{\boxspacing}
\begin{Verbatim}[commandchars=\\\{\}]
\PY{n}{path} \PY{o}{=} \PY{l+s+s2}{\PYZdq{}}\PY{l+s+s2}{C:/Users/ferqu/Documents/Mineria de datos/Caso\PYZsh{}2/hotel.csv}\PY{l+s+s2}{\PYZdq{}}
\PY{n}{ad} \PY{o}{=} \PY{n}{noSup}\PY{o}{.}\PY{n}{analisisEDA}\PY{p}{(}\PY{n}{path}\PY{p}{,} \PY{l+m+mi}{1}\PY{p}{)}
\PY{n}{df} \PY{o}{=} \PY{n}{ad}\PY{o}{.}\PY{n}{df}
\PY{n}{df}\PY{o}{.}\PY{n}{dtypes}
\end{Verbatim}
\end{tcolorbox}

            \begin{tcolorbox}[breakable, size=fbox, boxrule=.5pt, pad at break*=1mm, opacityfill=0]
\prompt{Out}{outcolor}{2}{\boxspacing}
\begin{Verbatim}[commandchars=\\\{\}]
hotel                              object
is\_canceled                         int64
lead\_time                           int64
arrival\_date\_month                 object
stays\_in\_weekend\_nights             int64
stays\_in\_week\_nights                int64
adults                              int64
children                            int64
babies                              int64
previous\_cancellations              int64
previous\_bookings\_not\_canceled      int64
assigned\_room\_type                 object
booking\_changes                     int64
deposit\_type                       object
days\_in\_waiting\_list                int64
customer\_type                      object
adr                               float64
required\_car\_parking\_spaces         int64
total\_of\_special\_requests           int64
reservation\_status                 object
dtype: object
\end{Verbatim}
\end{tcolorbox}
        
    \paragraph{ACP}\label{acp}

    \begin{tcolorbox}[breakable, size=fbox, boxrule=1pt, pad at break*=1mm,colback=cellbackground, colframe=cellborder]
\prompt{In}{incolor}{3}{\boxspacing}
\begin{Verbatim}[commandchars=\\\{\}]
\PY{n}{ad}\PY{o}{.}\PY{n}{analisisNumerico}\PY{p}{(}\PY{p}{)}
\PY{n}{acp} \PY{o}{=} \PY{n}{mf}\PY{o}{.}\PY{n}{NoSupervisado}\PY{p}{(}\PY{n}{ad}\PY{o}{.}\PY{n}{df}\PY{p}{)}
\PY{n}{acp}\PY{o}{.}\PY{n}{ACP}\PY{p}{(}\PY{n}{n\PYZus{}componentes}\PY{o}{=}\PY{l+m+mi}{2}\PY{p}{)}
\end{Verbatim}
\end{tcolorbox}

    \begin{center}
    \adjustimage{max size={0.9\linewidth}{0.9\paperheight}}{output_8_0.png}
    \end{center}
    { \hspace*{\fill} \\}
    
    \begin{center}
    \adjustimage{max size={0.9\linewidth}{0.9\paperheight}}{output_8_1.png}
    \end{center}
    { \hspace*{\fill} \\}
    
    \begin{center}
    \adjustimage{max size={0.9\linewidth}{0.9\paperheight}}{output_8_2.png}
    \end{center}
    { \hspace*{\fill} \\}
    
    \subparagraph{Observaciones:}\label{observaciones}

\begin{enumerate}
\def\labelenumi{\arabic{enumi}.}
\tightlist
\item
  Primera Imagen:
\end{enumerate}

\begin{itemize}
\tightlist
\item
  Los dos primeros componentes principales capturan juntos un 27.85\% de
  la varianza, lo cual te da una vision parcial pero util de la
  estructura del conjunto de datos original.
\item
  La mayoria de los puntos esta concentrada cerca del origen, indicando
  que comparten caracteristicas similares y posiblemente forman un grupo
  con patrones comunes.
\item
  Algunas etiquetas como 7733, 11259, 4071, 4452 y 4484 se alejan del
  centro, lo que sugiere que podrian ser outliers o casos con
  propiedades unicas.
\end{itemize}

\begin{enumerate}
\def\labelenumi{\arabic{enumi}.}
\setcounter{enumi}{1}
\tightlist
\item
  Segunda Imagen:
\end{enumerate}

\begin{itemize}
\tightlist
\item
  Las que apuntan en la misma direccion estan correlacionadas
  positivamente; las opuestas, negativamente.
\item
  Las variables con flechas largas (como adr o lead\_time) aportan mas a
  los ejes principales del ACP.
\item
  El eje 0 parece capturar aspectos de la reserva (duracion, precio),
  mientras que el eje 1 refleja comportamiento del cliente
  (cancelaciones, cambios).
\end{itemize}

\begin{enumerate}
\def\labelenumi{\arabic{enumi}.}
\setcounter{enumi}{2}
\tightlist
\item
  Tercera Imagen:
\end{enumerate}

\begin{itemize}
\tightlist
\item
  Observaciones alejadas indican casos con patrones unicos; podrian ser
  interesantes para analizar mas a fondo.
\item
  Agrupamiento central sugiere que la mayoria de las observaciones
  comparten caracteristicas similares.
\item
  Variables influyentes como lead\_time, adr y special\_requests tienen
  vectores largos, lo que indica que explican buena parte de la
  varianza.
\item
  Relaciones entre variables como stays\_in\_week\_nights y
  weekend\_nights (correlacion positiva), o lead\_time vs
  previous\_bookings\_not\_canceled (relacion negativa), ayudan a
  entender la estructura de los datos.
\end{itemize}

    \subsubsection{K-Medias}\label{k-medias}

    \begin{tcolorbox}[breakable, size=fbox, boxrule=1pt, pad at break*=1mm,colback=cellbackground, colframe=cellborder]
\prompt{In}{incolor}{4}{\boxspacing}
\begin{Verbatim}[commandchars=\\\{\}]
\PY{n}{ad}\PY{o}{.}\PY{n}{analisisNumerico}\PY{p}{(}\PY{p}{)}
\PY{n}{kmedias} \PY{o}{=} \PY{n}{mf}\PY{o}{.}\PY{n}{NoSupervisado}\PY{p}{(}\PY{n}{ad}\PY{o}{.}\PY{n}{df}\PY{p}{)}
\PY{n}{kmedias}\PY{o}{.}\PY{n}{KMedia}\PY{p}{(}\PY{p}{)}
\end{Verbatim}
\end{tcolorbox}

    \begin{center}
    \adjustimage{max size={0.9\linewidth}{0.9\paperheight}}{output_11_0.png}
    \end{center}
    { \hspace*{\fill} \\}
    
    \begin{center}
    \adjustimage{max size={0.9\linewidth}{0.9\paperheight}}{output_11_1.png}
    \end{center}
    { \hspace*{\fill} \\}
    
    \begin{center}
    \adjustimage{max size={0.9\linewidth}{0.9\paperheight}}{output_11_2.png}
    \end{center}
    { \hspace*{\fill} \\}
    
    \begin{Verbatim}[commandchars=\\\{\}]
c:\textbackslash{}Users\textbackslash{}ferqu\textbackslash{}Documents\textbackslash{}Mineria de datos\textbackslash{}Caso\#2\textbackslash{}M\_Caso1y2.py:295:
RuntimeWarning: invalid value encountered in divide
  max(n) != min(n) else (n/n * 50) for n in centros.T])
    \end{Verbatim}

    \begin{center}
    \adjustimage{max size={0.9\linewidth}{0.9\paperheight}}{output_11_4.png}
    \end{center}
    { \hspace*{\fill} \\}
    
    \paragraph{Observaciones}\label{observaciones}

Al analizar la grafica resultante del algoritmo K-medias, se identifica
una clara segmentacion de los datos en tres grupos bien definidos. La
proyeccion en dos componentes principales facilita la visualizacion de
la estructura interna del conjunto de datos. Se observa que los clusters
estan distribuidos de forma coherente, con minima superposicion, lo que
sugiere una diferenciacion significativa entre los grupos. Esta
representacion evidencia que el metodo logro captar patrones relevantes,
permitiendo una interpretacion mas precisa de la distribucion y
similitud de las observaciones.

Al examinar la distribucion de variables por cluster, se observa una
diferenciacion clara en las caracteristicas que definen a cada grupo. El
grafico permite identificar, por ejemplo, que el grupo 0 presenta
mayores valores en variables como ``adr'' y ``lead\_time'', lo cual
podria estar asociado a estadias mas planificadas y de mayor costo. Por
su parte, los clusters 1 y 2 destacan en variables como
``required\_car\_parking\_spaces'', lo que sugiere distintos perfiles de
necesidades logisticas. Esta representacion facilita la interpretacion
comparativa entre los grupos, permitiendo asociar patrones de
comportamiento o caracteristicas comunes dentro de cada cluster.

Al interpretar la grafica de radar, se evidencian diferencias marcadas
entre los tres clusters en relacion con diversas variables asociadas a
reservas hoteleras. Por ejemplo, el grupo 0 destaca por valores elevados
en ``lead\_time'' y ``adr'', lo que sugiere un perfil de reservas
anticipadas y con tarifas promedio mas altas. En contraste, el cluster 2
presenta mayor incidencia de cancelaciones previas y cambios en la
reserva, lo cual podria asociarse a un comportamiento mas inestable o
incierto. Por su parte, el cluster 1 muestra niveles mas altos en
variables como ``required\_car\_parking\_spaces'' y ``children'', lo que
podria reflejar un perfil mas familiar. Esta representacion permite
comparar de forma simultanea multiples atributos, brindando una vision
integral de las particularidades de cada grupo identificado por el
modelo de clustering.

    \subsubsection{HAC}\label{hac}

    \begin{tcolorbox}[breakable, size=fbox, boxrule=1pt, pad at break*=1mm,colback=cellbackground, colframe=cellborder]
\prompt{In}{incolor}{5}{\boxspacing}
\begin{Verbatim}[commandchars=\\\{\}]
\PY{n}{hac} \PY{o}{=} \PY{n}{mf}\PY{o}{.}\PY{n}{NoSupervisado}\PY{p}{(}\PY{n}{ad}\PY{o}{.}\PY{n}{df}\PY{p}{)}
\PY{n}{hac}\PY{o}{.}\PY{n}{HAC}\PY{p}{(}\PY{p}{)}
\end{Verbatim}
\end{tcolorbox}

    \begin{Verbatim}[commandchars=\\\{\}]
Metodo Ward:
    \end{Verbatim}

    \begin{center}
    \adjustimage{max size={0.9\linewidth}{0.9\paperheight}}{output_14_1.png}
    \end{center}
    { \hspace*{\fill} \\}
    
    \begin{Verbatim}[commandchars=\\\{\}]
Average:
    \end{Verbatim}

    \begin{center}
    \adjustimage{max size={0.9\linewidth}{0.9\paperheight}}{output_14_3.png}
    \end{center}
    { \hspace*{\fill} \\}
    
    \begin{Verbatim}[commandchars=\\\{\}]
Single:
    \end{Verbatim}

    \begin{center}
    \adjustimage{max size={0.9\linewidth}{0.9\paperheight}}{output_14_5.png}
    \end{center}
    { \hspace*{\fill} \\}
    
    \begin{Verbatim}[commandchars=\\\{\}]
Complete:
    \end{Verbatim}

    \begin{center}
    \adjustimage{max size={0.9\linewidth}{0.9\paperheight}}{output_14_7.png}
    \end{center}
    { \hspace*{\fill} \\}
    
    \begin{center}
    \adjustimage{max size={0.9\linewidth}{0.9\paperheight}}{output_14_8.png}
    \end{center}
    { \hspace*{\fill} \\}
    
    \begin{center}
    \adjustimage{max size={0.9\linewidth}{0.9\paperheight}}{output_14_9.png}
    \end{center}
    { \hspace*{\fill} \\}
    
    \begin{center}
    \adjustimage{max size={0.9\linewidth}{0.9\paperheight}}{output_14_10.png}
    \end{center}
    { \hspace*{\fill} \\}
    
    \begin{center}
    \adjustimage{max size={0.9\linewidth}{0.9\paperheight}}{output_14_11.png}
    \end{center}
    { \hspace*{\fill} \\}
    
    \begin{center}
    \adjustimage{max size={0.9\linewidth}{0.9\paperheight}}{output_14_12.png}
    \end{center}
    { \hspace*{\fill} \\}
    
    \paragraph{Observaciones:}\label{observaciones}

Se obtuvo una estructura dendrogramatica que refleja claramente como se
agrupan progresivamente las observaciones. El metodo de Ward busca
minimizar la varianza interna de los clusters en cada fusion, lo que
resulta en un agrupamiento mas compacto y homogeneo. En el dendrograma
se identifican tres grandes ramas principales, lo que sugiere una
segmentacion natural en tres clusters. Ademas, la distancia a la cual se
realizan las fusiones (representada por la altura de las uniones)
permite evaluar que tan distintos son los grupos antes de ser
combinados, destacando rupturas significativas que justifican la
eleccion del numero optimo de clusters. Esta visualizacion ofrece una
perspectiva jerarquica de la estructura del conjunto de datos y
complementa de manera efectiva los resultados obtenidos con metodos como
K-medias.

    \paragraph{T-SNE}\label{t-sne}

    \begin{tcolorbox}[breakable, size=fbox, boxrule=1pt, pad at break*=1mm,colback=cellbackground, colframe=cellborder]
\prompt{In}{incolor}{6}{\boxspacing}
\begin{Verbatim}[commandchars=\\\{\}]
\PY{n}{tsne} \PY{o}{=} \PY{n}{mf}\PY{o}{.}\PY{n}{NoSupervisado}\PY{p}{(}\PY{n}{ad}\PY{o}{.}\PY{n}{df}\PY{p}{)}
\PY{n}{tsne}\PY{o}{.}\PY{n}{TSNE}\PY{p}{(}\PY{p}{)}
\end{Verbatim}
\end{tcolorbox}

    \begin{Verbatim}[commandchars=\\\{\}]
Silhouette Score: 0.5550027
    \end{Verbatim}

    \begin{center}
    \adjustimage{max size={0.9\linewidth}{0.9\paperheight}}{output_17_1.png}
    \end{center}
    { \hspace*{\fill} \\}
    
    \begin{center}
    \adjustimage{max size={0.9\linewidth}{0.9\paperheight}}{output_17_2.png}
    \end{center}
    { \hspace*{\fill} \\}
    
    \paragraph{Observaciones}\label{observaciones}

Esto revela una clara correspondencia entre la estructura espacial de
las observaciones y las agrupaciones generadas. Se puede observar que
los grupos ocupan regiones bien diferenciadas dentro del plano
bidimensional, lo cual sugiere coherencia entre la segmentacion
efectuada por el modelo y las relaciones no lineales captadas por UMAP.
Esta visualizacion permite validar de forma visual la consistencia
interna de los clusters y aporta evidencia adicional sobre la presencia
de patrones naturales en los datos. Asimismo, refuerza la idea de que
las observaciones dentro de cada cluster comparten caracteristicas
similares, mientras que los grupos entre si son significativamente
distintos.

    \paragraph{UMAP}\label{umap}

    \begin{tcolorbox}[breakable, size=fbox, boxrule=1pt, pad at break*=1mm,colback=cellbackground, colframe=cellborder]
\prompt{In}{incolor}{7}{\boxspacing}
\begin{Verbatim}[commandchars=\\\{\}]
\PY{n}{umap} \PY{o}{=} \PY{n}{mf}\PY{o}{.}\PY{n}{NoSupervisado}\PY{p}{(}\PY{n}{ad}\PY{o}{.}\PY{n}{df}\PY{p}{)}
\PY{n}{umap}\PY{o}{.}\PY{n}{UMAP}\PY{p}{(}\PY{n}{n\PYZus{}componentes}\PY{o}{=}\PY{l+m+mi}{2}\PY{p}{,} \PY{n}{n\PYZus{}neighbors}\PY{o}{=}\PY{l+m+mi}{2}\PY{p}{)}
\end{Verbatim}
\end{tcolorbox}

    \begin{Verbatim}[commandchars=\\\{\}]
c:\textbackslash{}Users\textbackslash{}ferqu\textbackslash{}env\textbackslash{}Lib\textbackslash{}site-packages\textbackslash{}sklearn\textbackslash{}utils\textbackslash{}deprecation.py:151:
FutureWarning: 'force\_all\_finite' was renamed to 'ensure\_all\_finite' in 1.6 and
will be removed in 1.8.
  warnings.warn(
    \end{Verbatim}

    \begin{center}
    \adjustimage{max size={0.9\linewidth}{0.9\paperheight}}{output_20_1.png}
    \end{center}
    { \hspace*{\fill} \\}
    
            \begin{tcolorbox}[breakable, size=fbox, boxrule=.5pt, pad at break*=1mm, opacityfill=0]
\prompt{Out}{outcolor}{7}{\boxspacing}
\begin{Verbatim}[commandchars=\\\{\}]
array([[  5.5337305 ,  -0.993348  ],
       [  2.179117  ,  10.808589  ],
       [ -3.3342304 ,   0.7840228 ],
       [  6.771133  , -14.399791  ],
       [ -2.0037494 ,   9.948875  ],
       [ -1.4479944 ,   5.7463264 ],
       [  9.033091  ,  -2.8722503 ],
       [  2.1041133 ,  10.88014   ],
       [ -7.2219377 ,   2.0000758 ],
       [ 12.236504  ,  -3.6932182 ],
       [  6.468979  ,   2.842167  ],
       [  8.949531  ,   0.26713642],
       [ 11.806546  ,  -7.126201  ],
       [  2.850833  , -14.806117  ],
       [ 15.046131  ,  -8.770282  ],
       [  0.20914118,  -4.7024565 ],
       [ -3.3374803 ,   0.81150925],
       [ -1.4504361 ,   5.7522216 ],
       [  2.549704  ,   3.5305128 ],
       [  1.9342998 ,  -7.754767  ],
       [  6.586344  ,  -8.811409  ],
       [ -6.036551  ,   0.28534237],
       [  3.4622154 ,  -1.6255966 ],
       [  3.6541018 ,  -8.817361  ],
       [  7.946172  ,  -7.2407475 ],
       [  6.1423764 ,   5.514166  ],
       [ -4.171351  , -10.369179  ],
       [ 10.435369  ,  -3.588765  ],
       [  4.025318  ,   0.5759974 ],
       [ -2.6788108 , -13.099766  ],
       [  5.9988723 ,   5.3587213 ],
       [  1.8943089 , -11.033788  ],
       [-15.319459  ,   0.87237877],
       [ -5.1019273 ,  -7.788152  ],
       [-13.998784  ,  -9.678727  ],
       [ -2.4046059 , -13.379671  ],
       [ -2.7467105 , -13.03145   ],
       [-14.381382  ,   2.6937215 ],
       [ -2.5237563 , -13.256608  ],
       [  5.755248  ,  16.669693  ],
       [  0.8945923 ,  -2.9236908 ],
       [  1.0747726 ,  -5.2029386 ],
       [  2.6724415 , -14.637349  ],
       [  9.03339   ,  -2.8718927 ],
       [  1.0267948 ,   7.504248  ],
       [-11.36776   ,  11.210333  ],
       [  2.4645894 ,   5.369081  ],
       [  0.36693552,  10.764571  ],
       [ -4.1727295 , -10.37108   ],
       [ -4.2473526 ,  10.948168  ],
       [  8.054329  ,  -7.2850895 ],
       [  8.883787  , -11.693076  ],
       [  0.03869003,   4.2584705 ],
       [  6.8461456 ,  14.616017  ],
       [ 11.764221  ,   8.6246195 ],
       [ -0.40090328,  -8.3369    ],
       [  6.779143  ,  14.683615  ],
       [ -9.732835  ,   2.3545198 ],
       [-11.081235  ,   3.9076624 ],
       [  6.7811623 ,  14.680471  ],
       [ 10.435997  ,  -3.5891783 ],
       [  6.771157  , -14.39986   ],
       [  0.80606675,  -3.0069783 ],
       [-15.319448  ,   0.87237024],
       [  8.533703  ,  -7.0353675 ],
       [  2.6586823 , -14.623999  ],
       [  1.7083753 , -11.219349  ],
       [  5.832414  ,  16.74697   ],
       [ 17.381245  ,   3.1466277 ],
       [  8.883769  , -11.692986  ],
       [ -4.943009  , -11.161281  ],
       [ 12.236343  ,  -3.6937706 ],
       [  7.9499283 ,  -0.6294469 ],
       [ -5.9387507 ,   3.811874  ],
       [  0.87115514,   1.7152076 ],
       [ 14.975861  ,  -8.699718  ],
       [ -1.7615162 ,  -2.321627  ],
       [  5.929784  ,   5.283398  ],
       [ -2.8732111 ,  13.920083  ],
       [ -5.729157  ,   6.246438  ],
       [ -5.331322  ,  -3.0694792 ],
       [ 10.427133  ,  -3.5841293 ],
       [  5.534447  ,  -0.9948426 ],
       [ -5.030829  , -16.010239  ],
       [ -6.3504863 ,   9.156054  ],
       [  0.95488966,   7.57754   ],
       [ -4.155122  ,   6.0461607 ],
       [-10.461511  ,  -2.6786587 ],
       [ -4.1561294 ,   6.047151  ],
       [-10.30959   ,  -6.3040323 ],
       [ 17.394558  ,   3.1330671 ],
       [ -2.4663265 ,  -6.7935514 ],
       [  0.9738526 ,  13.08857   ],
       [ -2.507023  ,  12.746352  ],
       [ -4.24668   ,  10.948435  ],
       [  6.3129907 ,  -5.119085  ],
       [ -5.446764  ,  -1.3476611 ],
       [ 17.35539   ,   3.172678  ],
       [  6.4763393 ,  -4.956806  ],
       [ -9.875039  ,   5.830971  ],
       [ -7.06671   ,  -4.524791  ],
       [  2.5648968 ,   3.5539858 ],
       [  6.359189  ,  -5.0732136 ],
       [  0.6020865 ,   1.7890905 ],
       [  6.512512  ,   0.957706  ],
       [  6.415946  ,  -5.0175805 ],
       [-14.381639  ,   2.6935782 ],
       [-11.055443  ,   3.9328568 ],
       [ 11.883626  ,  -6.0922055 ],
       [  6.1013975 ,   5.4714084 ],
       [  3.662614  ,  -8.807647  ],
       [ -7.796108  ,  -2.1912827 ],
       [-12.441258  ,   4.369908  ],
       [  1.9342597 ,  -7.754897  ],
       [  5.2460237 ,   4.8637633 ],
       [-10.860102  ,   9.163494  ],
       [ -5.10519   , -15.935707  ],
       [  1.7293457 , -11.198892  ],
       [ -4.9865546 , -12.7703705 ],
       [ -4.9881907 , -16.052725  ],
       [  6.345929  ,  -2.7469487 ],
       [-11.055748  ,   3.9322248 ],
       [  2.5534358 ,   3.5471098 ],
       [ -4.9867826 , -12.769921  ],
       [  3.6043825 ,  -8.865922  ],
       [ -9.802118  ,   5.9142528 ],
       [ -0.400562  ,  -8.337136  ],
       [  2.292047  ,  10.695159  ],
       [  2.2974617 ,  10.688164  ],
       [  0.11119875,  -1.0979478 ],
       [ -7.434176  ,  -9.801782  ],
       [-10.859727  ,   9.163     ],
       [  1.8637997 , -11.064693  ],
       [ -9.261513  ,  -9.712748  ],
       [  0.25914475,  10.875382  ],
       [  8.967442  ,   0.2887511 ],
       [ 15.012135  ,  -8.736214  ],
       [  2.281969  ,   5.191028  ],
       [  6.484639  ,   0.9875201 ],
       [  7.9408402 ,  -0.6206865 ],
       [ -5.435779  ,  -1.3414458 ],
       [  0.03813242,   4.2592225 ],
       [  4.252962  ,   8.038654  ],
       [ -8.975315  ,  12.654426  ],
       [-12.4443655 ,   4.370999  ],
       [  0.1717122 ,  -1.0395496 ],
       [-10.277074  ,  -6.3387456 ],
       [  0.75699735,  -3.0523045 ],
       [ -4.274463  ,  -8.330859  ],
       [  8.554731  ,  -6.9726644 ],
       [ -4.278872  ,  -8.328975  ],
       [ 17.349398  ,   3.1786416 ],
       [ 12.40938   ,   1.4173777 ],
       [-10.860807  ,   9.162976  ],
       [  6.0570354 ,   5.429276  ],
       [ 12.400441  ,   1.4067891 ],
       [ -7.6605477 ,  -2.3328927 ],
       [ -2.0043921 ,   9.949544  ],
       [  2.3454473 ,   5.250984  ],
       [ 17.39148   ,   3.1361923 ],
       [  5.53469   ,  -0.99498105],
       [ -4.98532   , -12.769196  ],
       [ -4.1639524 , -10.361379  ],
       [  0.97366846,  13.088533  ],
       [  1.101089  ,  -5.226756  ],
       [  8.916844  ,   0.24048854],
       [ -3.1785583 ,   1.9810627 ],
       [  0.89132035,  -2.930569  ],
       [ 11.336409  ,   2.7478156 ],
       [  6.0293293 ,   5.4030514 ],
       [ -7.932261  ,  -3.9316664 ],
       [-11.94414   ,  -5.2604585 ],
       [  3.6654332 ,  -8.805857  ],
       [  7.93633   ,  -7.2285085 ],
       [  0.9548285 ,   0.96296966],
       [ -7.195493  ,   1.9748237 ],
       [  7.8875523 ,  -7.2465672 ],
       [ -0.3987897 ,  -8.335331  ],
       [ -0.40019232,  -8.337278  ],
       [ 12.398236  ,   1.4051822 ],
       [  2.583613  ,   3.5577948 ],
       [  1.8771398 , -11.050468  ],
       [  3.2929544 ,  -5.817016  ],
       [ -4.278021  ,  -8.327934  ],
       [ -7.437732  ,  -9.809508  ],
       [ -4.9569387 ,  -7.9330974 ],
       [  0.5108033 ,   1.8746519 ],
       [  5.2459745 ,   4.8637733 ],
       [  9.033569  ,  -2.8718977 ],
       [ -6.0435376 ,   0.27690598],
       [  2.5199757 ,   3.508791  ],
       [  1.0188022 ,   7.5128484 ],
       [  0.18424745,  -1.0286298 ],
       [  1.094495  ,  -5.2189417 ],
       [ -0.12202916,   8.127715  ],
       [ -7.605444  ,  -2.4067152 ],
       [ -5.2376223 ,  -2.980819  ],
       [ -7.999063  ,  -6.025628  ],
       [ -6.0782104 ,   0.24638942],
       [  6.49454   ,   2.8659952 ],
       [ -1.7490292 ,  -2.3330688 ],
       [  5.8219557 ,  16.736597  ],
       [ 11.782462  ,  -7.1493616 ],
       [  8.369578  ,  11.619588  ],
       [ -1.7722851 ,  -2.3088832 ],
       [ -7.999218  ,  -6.0244513 ],
       [  1.1045697 ,  -2.756105  ],
       [ -7.6955304 ,  -2.2980013 ],
       [  0.0873059 ,  -1.1216266 ],
       [ -7.1699495 ,   1.9458634 ],
       [ -2.8726878 ,  13.918443  ],
       [ 11.336939  ,   2.7472086 ],
       [  8.967542  ,   0.28209218],
       [ -7.9296584 ,  -3.9355667 ],
       [ 10.322921  ,  -0.32920167],
       [ -4.962128  , -16.078917  ],
       [ -5.086132  ,  -7.8046665 ],
       [ -8.95315   ,  -0.4405811 ],
       [  6.4762435 ,   2.8472507 ],
       [ 15.026671  ,   6.0647893 ],
       [ -7.0317163 ,  -4.5519896 ],
       [-11.367874  ,  11.2104845 ],
       [ 11.3345375 ,   2.7515655 ],
       [-10.460979  ,  -2.6783304 ],
       [  2.4422035 ,   5.3457355 ],
       [  0.22175145,  -4.712174  ],
       [  6.4509478 ,   1.0231183 ],
       [  2.2769918 ,  10.7104    ],
       [ -7.009197  ,  -4.573463  ],
       [-14.381813  ,   2.6934988 ],
       [  8.076585  ,  -7.2651215 ],
       [  7.275663  ,   7.991899  ],
       [-10.258357  ,  -6.3571105 ],
       [ -3.0460746 ,   1.8440179 ],
       [  3.4110332 ,  -1.5769163 ],
       [  6.52031   ,  10.182805  ],
       [ 11.82457   ,  -7.1078496 ],
       [  8.371735  ,  -7.1893272 ],
       [ 12.984816  ,   5.0706944 ],
       [ -3.336734  ,   0.7980301 ],
       [  0.35598963,  10.777889  ],
       [ -7.540916  ,  -2.4636023 ],
       [-10.461544  ,  -2.6784294 ],
       [ 12.318875  ,   1.3250282 ],
       [ -9.390702  ,  -7.083277  ],
       [ -5.3941684 ,  10.637261  ],
       [-12.291727  ,   0.607486  ],
       [ 12.362449  ,   1.3703527 ],
       [ -5.3014245 ,  -3.042563  ],
       [  5.5953083 ,   2.366217  ],
       [ -4.9362555 , -11.169212  ],
       [  6.1087832 ,   5.4791665 ],
       [ -3.2747438 ,  -0.80517423],
       [  3.356916  ,  -1.5216309 ],
       [ 11.764199  ,   8.62463   ],
       [ -0.18761098,   8.192601  ],
       [  2.6664445 , -14.63168   ],
       [  0.97371316,  13.088472  ],
       [  4.0258603 ,   0.5769582 ],
       [ -8.975282  ,  12.654433  ],
       [  4.25275   ,   8.038986  ],
       [  2.6525803 , -14.617477  ],
       [ -0.16283093,   8.168655  ],
       [ -5.4913216 ,  -3.218908  ],
       [  0.704932  ,   1.7116318 ],
       [ 15.064051  ,  -8.789977  ],
       [ -4.919554  ,   7.4528556 ],
       [  3.3617356 ,  -5.8855243 ],
       [ -5.0194664 ,  -7.870831  ],
       [  6.3099833 ,  -5.1219373 ],
       [-11.960116  ,  -5.245561  ],
       [ -5.617932  ,   6.3535924 ],
       [ 10.32237   ,  -0.330073  ],
       [  7.2754254 ,   7.9917645 ],
       [ 15.036426  ,  -8.760378  ],
       [ -6.0599923 ,   0.26796553],
       [ -6.3500423 ,   9.156553  ],
       [ -5.660585  ,   6.3133082 ],
       [  3.228039  ,  -5.7526574 ],
       [  0.16130292,  -4.651835  ],
       [  0.99215513,   7.5387626 ],
       [  8.36957   ,  11.619453  ],
       [  5.8116856 ,   5.1612463 ],
       [  5.5345955 ,  -0.99514234],
       [  5.6128864 ,  16.526148  ],
       [ -2.565356  , -13.212766  ],
       [ -1.6396589 ,   8.194168  ],
       [ -9.732357  ,   2.3546858 ],
       [  3.2630928 ,  -5.787542  ],
       [ -4.1540093 ,   6.045831  ],
       [ -2.686734  , -13.090883  ],
       [  7.275833  ,   7.9914584 ],
       [ -4.140892  , -10.339344  ],
       [ -2.3976588 ,  -5.371046  ],
       [ -3.1505494 ,   1.9520545 ],
       [  8.507753  ,  -7.1049733 ],
       [  2.1266212 ,  10.858533  ],
       [  3.2688844 ,  -1.4309406 ],
       [ -2.321645  , -13.462577  ],
       [ -2.4197676 ,  -5.3949637 ],
       [  2.187727  ,   5.09727   ],
       [  7.927003  ,  -0.6068985 ],
       [ -2.4607103 , -13.322474  ],
       [  6.8147883 ,  14.647084  ],
       [  0.9310428 ,   0.9797541 ],
       [ -1.8938279 ,   4.0671453 ],
       [ -8.883836  ,  -0.37147534],
       [ -7.7458243 ,  -2.242368  ],
       [ -2.4460483 ,  -5.422641  ],
       [  0.96042114,  -2.8818612 ],
       [ -9.417407  ,   9.3531    ],
       [ -3.2780323 ,  -0.805858  ],
       [ -1.8934131 ,   4.0668035 ],
       [ -5.9387097 ,   3.8118143 ],
       [  0.8801714 ,  -2.9376554 ],
       [  6.520198  ,  10.183179  ],
       [ -2.4661357 ,  -6.7934504 ],
       [ -9.390438  ,  -7.0873313 ],
       [ 15.02661   ,   6.0646663 ],
       [ -3.3339193 ,   0.77905095],
       [  8.507934  ,  -6.9925575 ],
       [ 11.883618  ,  -6.092185  ],
       [ -9.420681  ,   9.352363  ],
       [ -7.058487  ,  -4.5295787 ],
       [  1.162419  ,  -2.704887  ],
       [ -5.4498277 ,  -3.1839304 ],
       [  4.252118  ,   8.039392  ],
       [ -7.4347057 ,  -9.803596  ],
       [  7.937945  ,  -0.6167733 ],
       [ -3.0062566 ,   1.7967815 ],
       [ -5.437234  ,  -1.3427477 ],
       [ -2.4292688 ,  -5.4078517 ],
       [-12.015736  ,  -5.1892447 ],
       [ -2.4288392 ,  -5.407291  ],
       [  3.602173  ,  -8.868784  ],
       [-12.282446  ,   0.6036096 ],
       [ -5.7308435 ,   6.2487    ],
       [ -9.39028   ,  -7.090733  ],
       [ -5.3954887 ,  10.6370125 ],
       [ -2.5070362 ,  12.746422  ],
       [ -5.3948574 ,  10.637556  ],
       [ -1.6394788 ,   8.194365  ],
       [  3.2298675 ,  -1.3883915 ],
       [  5.550538  ,  16.46225   ],
       [ -2.9456918 ,   1.7388344 ],
       [ -1.6978343 ,  -2.3847804 ],
       [  1.9122975 , -11.016042  ],
       [-11.205716  ,   3.784398  ],
       [  3.1908367 ,  -1.3474158 ],
       [  1.6904023 , -11.238354  ],
       [  6.3439417 ,  -2.7412302 ],
       [ -8.914151  ,  -0.3998333 ],
       [  6.4393797 ,   1.0359291 ],
       [ -4.15152   , -10.349462  ],
       [  6.4559865 ,   2.8260827 ],
       [  2.3846452 ,   5.2879133 ],
       [  6.520452  ,  10.183131  ],
       [ -2.8843713 ,   1.6781965 ],
       [  0.526176  ,   1.8600563 ],
       [ -5.5167103 ,  -3.2404778 ],
       [ -1.7245055 ,  -2.3576345 ],
       [ -7.5067916 ,  -2.48951   ],
       [ -9.2471485 ,  -9.698684  ],
       [ -8.867069  ,  -0.35444883],
       [ 11.764201  ,   8.624606  ],
       [ -5.1058517 ,  -7.7860394 ],
       [ -4.9791346 , -11.12697   ],
       [  5.5924797 ,   2.3637764 ],
       [ -9.293425  ,  -9.746106  ],
       [ -4.179067  ,   6.059521  ],
       [ 15.016527  ,  -8.740781  ],
       [  5.64394   ,  16.556389  ],
       [ -5.395362  ,  10.636842  ],
       [ -9.812534  ,   5.901845  ],
       [-13.998763  ,  -9.6787195 ],
       [  6.58694   ,  -8.81088   ],
       [-12.413192  ,   4.3584256 ],
       [  1.0221769 ,   7.509483  ],
       [  2.8381002 , -14.800958  ],
       [  1.1082206 ,  -5.2359643 ],
       [ 11.788646  ,  -7.143214  ],
       [ 12.236513  ,  -3.6933138 ],
       [-10.530316  ,  -6.08344   ],
       [-11.997216  ,  -5.2086706 ],
       [ -0.2009821 ,   8.206053  ],
       [-10.518907  ,  -6.0954647 ],
       [  0.82107216,   1.6705147 ],
       [ -3.3366177 ,   0.81520396],
       [ -9.420266  ,   9.352416  ],
       [ -4.9966497 , -16.045639  ],
       [-10.858497  ,   9.164016  ],
       [-11.367823  ,  11.210442  ],
       [-10.44264   ,  -6.1715207 ],
       [  0.31170714,  10.823095  ],
       [ -7.178097  ,   1.9545825 ],
       [  8.3973675 ,  -7.1859255 ],
       [ -7.1951923 ,   1.9740232 ],
       [ -5.7293777 ,   6.249786  ],
       [  6.771164  , -14.3998375 ],
       [  1.7368122 , -11.1917305 ],
       [ -5.0403657 , -16.000679  ],
       [-11.231418  ,   3.7549791 ],
       [  2.5770624 ,   3.537925  ],
       [ -4.9202724 ,   7.4517207 ],
       [ 12.984804  ,   5.0707793 ],
       [  0.5260206 ,   1.8611386 ],
       [-12.288002  ,   0.60569876],
       [  0.25332037,  10.8812475 ]], dtype=float32)
\end{Verbatim}
\end{tcolorbox}
        
    \paragraph{Observaciones}\label{observaciones}

La representacion generada mediante UMAP permite visualizar la
estructura interna del conjunto de datos proyectado en dos dimensiones.
A traves de esta tecnica de reduccion no lineal, se identifican
concentraciones y separaciones entre observaciones que podrian
corresponder a patrones o grupos latentes en los datos originales.
Aunque los puntos aparecen sin etiquetar o colorear por cluster en esta
grafica, la dispersion y densidad relativa de las regiones sugiere la
posible existencia de agrupamientos naturales. Esta visualizacion
resulta especialmente util para complementar metodos como clustering
jerarquico o K-medias, ya que permite validar visualmente la coherencia
espacial de las segmentaciones propuestas.

    \subsection{Dataset 2: Banco}\label{dataset-2-banco}

    \begin{tcolorbox}[breakable, size=fbox, boxrule=1pt, pad at break*=1mm,colback=cellbackground, colframe=cellborder]
\prompt{In}{incolor}{8}{\boxspacing}
\begin{Verbatim}[commandchars=\\\{\}]
\PY{n}{path} \PY{o}{=} \PY{l+s+s2}{\PYZdq{}}\PY{l+s+s2}{C:/Users/ferqu/Documents/Mineria de datos/Caso\PYZsh{}2/bank.csv}\PY{l+s+s2}{\PYZdq{}}
\PY{n}{ab} \PY{o}{=} \PY{n}{mf}\PY{o}{.}\PY{n}{analisisEDA}\PY{p}{(}\PY{n}{path}\PY{p}{,} \PY{l+m+mi}{1}\PY{p}{)}
\PY{n}{df} \PY{o}{=} \PY{n}{ab}\PY{o}{.}\PY{n}{df}
\end{Verbatim}
\end{tcolorbox}

    \subsubsection{ACP}\label{acp}

    \begin{tcolorbox}[breakable, size=fbox, boxrule=1pt, pad at break*=1mm,colback=cellbackground, colframe=cellborder]
\prompt{In}{incolor}{9}{\boxspacing}
\begin{Verbatim}[commandchars=\\\{\}]
\PY{n}{ab}\PY{o}{.}\PY{n}{analisisNumerico}\PY{p}{(}\PY{p}{)}
\PY{n}{acp} \PY{o}{=} \PY{n}{mf}\PY{o}{.}\PY{n}{NoSupervisado}\PY{p}{(}\PY{n}{ab}\PY{o}{.}\PY{n}{df}\PY{p}{)}
\PY{n}{acp}\PY{o}{.}\PY{n}{ACP}\PY{p}{(}\PY{n}{n\PYZus{}componentes}\PY{o}{=}\PY{l+m+mi}{2}\PY{p}{)}
\end{Verbatim}
\end{tcolorbox}

    \begin{center}
    \adjustimage{max size={0.9\linewidth}{0.9\paperheight}}{output_25_0.png}
    \end{center}
    { \hspace*{\fill} \\}
    
    \begin{center}
    \adjustimage{max size={0.9\linewidth}{0.9\paperheight}}{output_25_1.png}
    \end{center}
    { \hspace*{\fill} \\}
    
    \begin{center}
    \adjustimage{max size={0.9\linewidth}{0.9\paperheight}}{output_25_2.png}
    \end{center}
    { \hspace*{\fill} \\}
    
    \subsubsection{K-MEDIAS}\label{k-medias}

    \begin{tcolorbox}[breakable, size=fbox, boxrule=1pt, pad at break*=1mm,colback=cellbackground, colframe=cellborder]
\prompt{In}{incolor}{10}{\boxspacing}
\begin{Verbatim}[commandchars=\\\{\}]
\PY{n}{ab}\PY{o}{.}\PY{n}{analisisNumerico}\PY{p}{(}\PY{p}{)}
\PY{n}{kmedias} \PY{o}{=} \PY{n}{mf}\PY{o}{.}\PY{n}{NoSupervisado}\PY{p}{(}\PY{n}{ab}\PY{o}{.}\PY{n}{df}\PY{p}{)}
\PY{n}{kmedias}\PY{o}{.}\PY{n}{KMedia}\PY{p}{(}\PY{p}{)}
\end{Verbatim}
\end{tcolorbox}

    \begin{center}
    \adjustimage{max size={0.9\linewidth}{0.9\paperheight}}{output_27_0.png}
    \end{center}
    { \hspace*{\fill} \\}
    
    \begin{center}
    \adjustimage{max size={0.9\linewidth}{0.9\paperheight}}{output_27_1.png}
    \end{center}
    { \hspace*{\fill} \\}
    
    \begin{center}
    \adjustimage{max size={0.9\linewidth}{0.9\paperheight}}{output_27_2.png}
    \end{center}
    { \hspace*{\fill} \\}
    
    \begin{center}
    \adjustimage{max size={0.9\linewidth}{0.9\paperheight}}{output_27_3.png}
    \end{center}
    { \hspace*{\fill} \\}
    
    \subsubsection{HAC}\label{hac}

    \begin{tcolorbox}[breakable, size=fbox, boxrule=1pt, pad at break*=1mm,colback=cellbackground, colframe=cellborder]
\prompt{In}{incolor}{11}{\boxspacing}
\begin{Verbatim}[commandchars=\\\{\}]
\PY{n}{hac} \PY{o}{=} \PY{n}{mf}\PY{o}{.}\PY{n}{NoSupervisado}\PY{p}{(}\PY{n}{ab}\PY{o}{.}\PY{n}{df}\PY{p}{)}
\PY{n}{hac}\PY{o}{.}\PY{n}{HAC}\PY{p}{(}\PY{p}{)}
\end{Verbatim}
\end{tcolorbox}

    \begin{Verbatim}[commandchars=\\\{\}]
Metodo Ward:
    \end{Verbatim}

    \begin{center}
    \adjustimage{max size={0.9\linewidth}{0.9\paperheight}}{output_29_1.png}
    \end{center}
    { \hspace*{\fill} \\}
    
    \begin{Verbatim}[commandchars=\\\{\}]
Average:
    \end{Verbatim}

    \begin{center}
    \adjustimage{max size={0.9\linewidth}{0.9\paperheight}}{output_29_3.png}
    \end{center}
    { \hspace*{\fill} \\}
    
    \begin{Verbatim}[commandchars=\\\{\}]
Single:
    \end{Verbatim}

    \begin{center}
    \adjustimage{max size={0.9\linewidth}{0.9\paperheight}}{output_29_5.png}
    \end{center}
    { \hspace*{\fill} \\}
    
    \begin{Verbatim}[commandchars=\\\{\}]
Complete:
    \end{Verbatim}

    \begin{center}
    \adjustimage{max size={0.9\linewidth}{0.9\paperheight}}{output_29_7.png}
    \end{center}
    { \hspace*{\fill} \\}
    
    \begin{center}
    \adjustimage{max size={0.9\linewidth}{0.9\paperheight}}{output_29_8.png}
    \end{center}
    { \hspace*{\fill} \\}
    
    \begin{center}
    \adjustimage{max size={0.9\linewidth}{0.9\paperheight}}{output_29_9.png}
    \end{center}
    { \hspace*{\fill} \\}
    
    \begin{center}
    \adjustimage{max size={0.9\linewidth}{0.9\paperheight}}{output_29_10.png}
    \end{center}
    { \hspace*{\fill} \\}
    
    \begin{center}
    \adjustimage{max size={0.9\linewidth}{0.9\paperheight}}{output_29_11.png}
    \end{center}
    { \hspace*{\fill} \\}
    
    \begin{center}
    \adjustimage{max size={0.9\linewidth}{0.9\paperheight}}{output_29_12.png}
    \end{center}
    { \hspace*{\fill} \\}
    
    \subsubsection{TSNE}\label{tsne}

    \begin{tcolorbox}[breakable, size=fbox, boxrule=1pt, pad at break*=1mm,colback=cellbackground, colframe=cellborder]
\prompt{In}{incolor}{12}{\boxspacing}
\begin{Verbatim}[commandchars=\\\{\}]
\PY{n}{tsne} \PY{o}{=} \PY{n}{mf}\PY{o}{.}\PY{n}{NoSupervisado}\PY{p}{(}\PY{n}{ab}\PY{o}{.}\PY{n}{df}\PY{p}{)}
\PY{n}{tsne}\PY{o}{.}\PY{n}{TSNE}\PY{p}{(}\PY{p}{)}
\end{Verbatim}
\end{tcolorbox}

    \begin{Verbatim}[commandchars=\\\{\}]
Silhouette Score: 0.51973784
    \end{Verbatim}

    \begin{center}
    \adjustimage{max size={0.9\linewidth}{0.9\paperheight}}{output_31_1.png}
    \end{center}
    { \hspace*{\fill} \\}
    
    \begin{center}
    \adjustimage{max size={0.9\linewidth}{0.9\paperheight}}{output_31_2.png}
    \end{center}
    { \hspace*{\fill} \\}
    
    \subsubsection{UMAP}\label{umap}

    \begin{tcolorbox}[breakable, size=fbox, boxrule=1pt, pad at break*=1mm,colback=cellbackground, colframe=cellborder]
\prompt{In}{incolor}{13}{\boxspacing}
\begin{Verbatim}[commandchars=\\\{\}]
\PY{n}{umap} \PY{o}{=} \PY{n}{mf}\PY{o}{.}\PY{n}{NoSupervisado}\PY{p}{(}\PY{n}{ab}\PY{o}{.}\PY{n}{df}\PY{p}{)}
\PY{n}{umap}\PY{o}{.}\PY{n}{UMAP}\PY{p}{(}\PY{n}{n\PYZus{}componentes}\PY{o}{=}\PY{l+m+mi}{2}\PY{p}{,} \PY{n}{n\PYZus{}neighbors}\PY{o}{=}\PY{l+m+mi}{2}\PY{p}{)}
\end{Verbatim}
\end{tcolorbox}

    \begin{Verbatim}[commandchars=\\\{\}]
c:\textbackslash{}Users\textbackslash{}ferqu\textbackslash{}env\textbackslash{}Lib\textbackslash{}site-packages\textbackslash{}sklearn\textbackslash{}utils\textbackslash{}deprecation.py:151:
FutureWarning: 'force\_all\_finite' was renamed to 'ensure\_all\_finite' in 1.6 and
will be removed in 1.8.
  warnings.warn(
    \end{Verbatim}

    \begin{center}
    \adjustimage{max size={0.9\linewidth}{0.9\paperheight}}{output_33_1.png}
    \end{center}
    { \hspace*{\fill} \\}
    
            \begin{tcolorbox}[breakable, size=fbox, boxrule=.5pt, pad at break*=1mm, opacityfill=0]
\prompt{Out}{outcolor}{13}{\boxspacing}
\begin{Verbatim}[commandchars=\\\{\}]
array([[-4.219841 , -1.3147312],
       [ 2.6727483,  7.0485916],
       [-6.390315 , -1.4340844],
       {\ldots},
       [10.513667 , -3.1896482],
       [-4.1414857, -6.121375 ],
       [ 4.027965 ,  5.94341  ]], dtype=float32)
\end{Verbatim}
\end{tcolorbox}
        

    % Add a bibliography block to the postdoc
    
    
    
\end{document}
